% =============================================================================
%  Chapter 9 — ElementGeometry Refactor
% =============================================================================

\chapter{ElementGeometry Refactor}
\label{ch:element_geometry_refactor}

This chapter documents the refactor of the \texttt{ElementGeometry} subsystem.
The goal of the change was not to redesign the whole geometry layer from
scratch, but to fix the most important semantic and architectural weaknesses
without degrading the hot path of element integration.

The guiding constraint was explicit:
\emph{prefer compile-time cost over runtime cost whenever the trade-off is
reasonable}.  In other words, stronger type constraints, cleaner ownership,
and slightly more template work at compile time are acceptable, whereas extra
runtime indirection inside the Gauss loop is not.


% ─────────────────────────────────────────────────────────────────────────
\section{Problem Statement}
\label{sec:eg_problem_statement}

Before the refactor, \texttt{ElementGeometry<dim>} already had a good high-level
idea: concrete geometry and quadrature strategy were stored behind a
type-erased, value-semantic wrapper.  However, three problems were
architecturally significant.

\subsection{Broken value semantics}

The wrapper cloned only its internal \texttt{pimpl\_}, while metadata such as
\texttt{integration\_point\_}, \texttt{sieve\_id}, and
\texttt{physical\_group\_} lived outside the cloned state.  As a consequence,
copy construction and copy assignment were semantically incomplete: two
\texttt{ElementGeometry} objects could share the same geometric map while
disagreeing on already-materialised Gauss-point data or mesh metadata.

This violated a core expectation of value types:
\[
\texttt{b = a} \quad \Longrightarrow \quad \texttt{b} \text{ must represent the same logical value as } \texttt{a}.
\]

\subsection{Integration strategy mixed into concrete geometry headers}

The concrete geometry headers contained both the element type and the
quadrature strategy implementation.  This worked, but it blurred the boundary
between two distinct concepts:

\begin{itemize}[nosep]
  \item a geometric interpolation on a reference cell, and
  \item a numerical integration policy on that reference cell.
\end{itemize}

Architecturally this was undesirable because the element family and the
quadrature family evolve for different reasons.

\subsection{Weak compile-time invariants}

Some invalid states were still representable:

\begin{itemize}[nosep]
  \item unsupported Gauss--Legendre orders silently produced zero-filled arrays;
  \item \texttt{detJ()} was visible even for embedded geometries with
        non-square Jacobians;
  \item tensor-product geometries did not consistently enforce
        \(\texttt{topological\_dim} \le \texttt{dim}\).
\end{itemize}

These were not merely cosmetic issues.  They weakened the library contract and
made failures harder to interpret.


% ─────────────────────────────────────────────────────────────────────────
\section{Design Goals}
\label{sec:eg_design_goals}

The refactor was driven by the following goals:

\begin{enumerate}[nosep]
  \item Preserve the existing \emph{external polymorphism} structure.
  \item Keep quadrature evaluation monomorphic and inlineable inside the
        concrete model.
  \item Repair value semantics of the wrapper.
  \item Make invalid geometric or quadrature configurations fail early.
  \item Reduce conceptual coupling between geometry and integration.
\end{enumerate}

Equally important is what was \emph{not} attempted in this refactor:
the geometry layer still depends on \texttt{Node<dim>}, which means it is not
yet a purely geometric vertex layer.  That larger ontological cleanup remains a
future step because it cuts across \texttt{Model}, DOF storage, and multiple
element families.


% ─────────────────────────────────────────────────────────────────────────
\section{Refactor Overview}
\label{sec:eg_refactor_overview}

\subsection{Metadata moved under wrapper ownership}

\texttt{ElementGeometry} now owns an internal metadata aggregate containing:

\begin{itemize}[nosep]
  \item the materialised integration points,
  \item the optional DMPlex sieve identifier, and
  \item the physical-group label imported from Gmsh.
\end{itemize}

This has two consequences.

First, copy construction and copy assignment now copy both the cloned geometry
implementation and the metadata, restoring proper value semantics.

Second, the metadata is no longer a public bag of fields.  It is accessed
through a narrow API:
\texttt{sieve\_id()}, \texttt{has\_sieve\_id()},
\texttt{physical\_group()}, \texttt{integration\_points()}, etc.

This is a direct SRP improvement: the wrapper remains responsible for its own
logical state, and downstream code no longer reaches into its internals.

\subsection{Quadrature strategies extracted to dedicated headers}

Two strategy types were extracted:

\begin{itemize}[nosep]
  \item \texttt{GaussLegendreCellIntegrator.hh}
  \item \texttt{SimplexIntegrator.hh}
\end{itemize}

The concrete geometry classes remain pure descriptions of interpolation and
mapping.  The numerical integration policy is still attached at compile time,
but it is no longer embedded in the same header for incidental reasons.

This separation is conceptually important.  In isoparametric finite elements,
geometry defines the map
\[
\mathbf{x}(\boldsymbol{\xi}) = \sum_{a=1}^{n_n} N_a(\boldsymbol{\xi})\,\mathbf{x}_a,
\]
while quadrature defines an approximation of
\[
\int_{\hat{\Omega}} f(\boldsymbol{\xi}) \,\mathrm{d}\hat{\Omega}
\;\approx\;
\sum_{g=1}^{n_g} w_g\,f(\boldsymbol{\xi}_g).
\]
They are related in use, but they are not the same abstraction.

\subsection{A small quadrature concept}

The new \texttt{QuadratureStrategy} concept formalises the contract expected by
\texttt{ElementGeometry}'s internal model:

\begin{itemize}[nosep]
  \item \texttt{num\_integration\_points}
  \item \texttt{reference\_integration\_point(i)}
  \item \texttt{weight(i)}
\end{itemize}

The benefit is not runtime polymorphism.  The benefit is that failures move to
compile time, where the error is cheaper and more local.


% ─────────────────────────────────────────────────────────────────────────
\section{Geometry and Measure}
\label{sec:eg_geometry_measure}

The refactor keeps the existing zero-cost logic for the differential measure.
Given the Jacobian \(J = \partial \mathbf{x} / \partial \boldsymbol{\xi}\), the
integration measure depends on the relationship between topological and
embedding dimensions:

\begin{align}
\text{full-dimensional cell:} \quad
& \mathrm{d}\Omega = |\det J|\,\mathrm{d}\hat{\Omega}, \\
\text{embedded curve:} \quad
& \mathrm{d}s = \lVert J \rVert\,\mathrm{d}\hat{s}, \\
\text{embedded surface in } \mathbb{R}^3: \quad
& \mathrm{d}S = \lVert J_1 \times J_2 \rVert\,\mathrm{d}\hat{S}.
\end{align}

This dispatch still happens inside
\texttt{OwningModel\_ElementGeometry<ElementType, IntegrationStrategy>} using
\texttt{if constexpr}.  Therefore the compiler sees the concrete element type,
the concrete topological dimension, and the concrete matrix shape.  No extra
runtime branch was introduced inside the integration loop.


% ─────────────────────────────────────────────────────────────────────────
\section{API Hardening}
\label{sec:eg_api_hardening}

\subsection{Restricting \texttt{detJ()}}

\texttt{detJ()} now exists only when the Jacobian is square, namely when the
element is full-dimensional:
\[
\texttt{topological\_dim} = \texttt{dim}.
\]

This is important because a determinant is not defined for the Jacobian of an
embedded beam or surface.  Exposing the member unconditionally was misleading:
it suggested a mathematically meaningful operation where none exists.

\subsection{Fail-fast Gauss--Legendre tables}

Unsupported Gauss--Legendre orders no longer fall back to zero arrays.
Instead, they trigger a compile-time error.  This is the right contract for a
template-based numerical kernel: if the requested order is unsupported, the
program should fail where the invalid instantiation occurs.

\subsection{Safer tensor-product constraints}

\texttt{LagrangeElement<Dim, N...>} now enforces that the number of natural
coordinates does not exceed the embedding dimension.  This prevents invalid
type combinations from entering the system.


% ─────────────────────────────────────────────────────────────────────────
\section{Materialised Integration Points}
\label{sec:eg_integration_points}

The wrapper continues to precompute physical Gauss-point coordinates when the
domain requests them.  The refactor extends this step so that the integration
weight is also stored in each materialised point.

This is useful for post-processing and point-wise objects attached to Gauss
points because the point now carries a coherent local state:

\begin{itemize}[nosep]
  \item physical coordinates,
  \item global identifier for VTK export,
  \item quadrature weight.
\end{itemize}

The operation remains amortised and explicit.  It does not change the
quadrature path used in stiffness or residual evaluation.


% ─────────────────────────────────────────────────────────────────────────
\section{Performance Discussion}
\label{sec:eg_performance}

The refactor was intentionally biased toward preserving runtime behaviour.

\subsection{What did not change}

\begin{itemize}[nosep]
  \item Concrete geometry evaluation is still template-based.
  \item The Jacobian and shape-function evaluation remain statically resolved.
  \item The quadrature loop still uses direct calls on the concrete strategy
        stored inside the model.
  \item No new virtual call was inserted inside \texttt{integrate()}.
\end{itemize}

\subsection{What did change}

\begin{itemize}[nosep]
  \item More errors are now diagnosed at compile time.
  \item Integration strategies live in dedicated headers, reducing conceptual
        coupling and making future dependency management easier.
  \item Copying \texttt{ElementGeometry} is slightly more expensive if
        integration points have already been materialised, because the copy now
        correctly copies them.  This is not a regression; it is the real cost
        of correct value semantics.
\end{itemize}

This is the intended trade-off.  Copying a fully materialised geometric entity
is not a hot-path operation in FEM assembly, whereas incorrect copies are
architecturally dangerous.


% ─────────────────────────────────────────────────────────────────────────
\section{What Was Deliberately Deferred}
\label{sec:eg_deferred}

The refactor does \emph{not} yet solve the deeper issue that the geometry layer
depends on \texttt{Node<dim>}, and \texttt{Node<dim>} already contains DOF
state.  From a strict modelling standpoint, a geometric layer should ideally
depend on a lighter vertex concept carrying only:

\begin{itemize}[nosep]
  \item identifier,
  \item coordinates,
  \item optional mesh-backend handle.
\end{itemize}

That separation would improve the ontology of the codebase and further align
the design with SRP and DIP.  However, it affects much more than
\texttt{ElementGeometry}; it is therefore better handled as a dedicated
follow-up refactor instead of being folded into the present change.


% ─────────────────────────────────────────────────────────────────────────
\section{Result}
\label{sec:eg_result}

Architecturally, the refactor moves \texttt{ElementGeometry} from
``good idea with leaky semantics'' toward ``stable boundary abstraction''.

The most important outcomes are:

\begin{enumerate}[nosep]
  \item wrapper copies are now semantically correct;
  \item geometry and quadrature are more cleanly separated;
  \item invalid numerical configurations fail early;
  \item the hot numerical path remains compile-time resolved.
\end{enumerate}

This is a meaningful improvement even though the subsystem is not yet the final
form of the geometry layer.  It removes a real correctness bug, strengthens the
public contract, and prepares the code for a later, larger cleanup around
geometric vertices versus analysis nodes.
